\documentclass[11pt,a4paper]{article}
\usepackage{jheppub}
\usepackage{booktabs}
\usepackage{siunitx}
\usepackage{microtype}
\hypersetup{pdftitle={Momentum-growing neutrino-quark interactions and the KM3NeT event: v0.9 audit},pdfauthor={Subhendu Rakshit}}

\title{Momentum-growing neutrino--quark interactions and the KM3NeT event:\\
an analytic and public-data audit (v0.9)}
\author{Subhendu Rakshit}
\abstract{We revisit a dimension-seven momentum-growing neutrino--quark
interaction after the observation of KM3-230213A and the 12.6-year IceCube
extremely-high-energy (EHE) search.  A leading-order parton calculation with
the conservative partonic-energy requirement $\sqrt{\hat s}<\Lambda$ retains
order-one cross-section changes at hundreds of PeV for a collider-bound
coefficient.  We derive an exact common-flux normalization identity and
audit the detector topology: the Case-6 operator has a neutral final state,
whereas the published KM3NeT response selects a bright muon track.  We
validate the released KM3NeT Ehlert-spectrum benchmark to 0.006\% and use
IceCube's released differential-limit curve in its stated sliding one-decade
$E^{-1}$ windows.  In an explicitly labelled all-flavor response proxy, the
limiting common-flux scales are 0.155, 0.128, 0.102 and 0.0747 for
$\Lambda=1.5,2.0,2.44,3.0$ TeV.  The corresponding unmodified direct
bright-track expectations are 0.0376, 0.0311, 0.0249 and 0.0181 events.
These are screening diagnostics, not frequentist 90\% limits: the public
release lacks the reconstructed-bin information and beyond-Standard-Model
response needed to reproduce the published likelihood.  We identify the
required topology and transport calculations for a definitive test and
state precise conditions under which alternative interactions could evade
the common-flux suppression.}
\keywords{Neutrino interactions, effective field theory, cosmic neutrinos,
IceCube, KM3NeT}

\begin{document}
\maketitle
\flushbottom
\tableofcontents

\section{Introduction and scope}
The KM3NeT collaboration reported a very bright track, KM3-230213A, made
by a muon traversing the partially built ARCA array.  For an $E_\nu^{-2}$
incident spectrum its median parent-neutrino energy is about $220$ PeV and
its stated central 90\% range is approximately $72$ PeV--$2.6$ EeV
\cite{KM3NeT:2025Nature}.  The event invites a quantitative comparison with
IceCube's search for neutrinos above a few PeV, which found three PeV-scale
events in 12.6 years but no event well above $10$ PeV
\cite{IceCube:2025EHE}.  The question here is narrower than explaining the
origin of one event: can the momentum growth of a laboratory-constrained
neutrino--quark operator increase an ARCA bright-track expectation after the
same flux passes the IceCube test?

We use an Ehlert-type cosmogenic spectrum as a fixed shape, with a single
free normalization $\kappa$ \cite{Ehlert:2024}.  The leading-order
cross-section table and the published detector effective areas are kept
separate.  This makes the structural result independent of the numerical
integral, while exposing the point where topology and response information
are missing.  Version 0.9 is a reproducible public-data audit.  It does not
claim to implement the collaboration's event-level EHE likelihood.

\section{Literature and novelty}
\subsection{Earth transport and neutrino deep-inelastic scattering}
Tau-neutrino propagation differs from pure absorption because charged-current
production of a tau followed by decay repopulates the neutrino flux at lower
energies.  Dutta, Reno and Sarcevic developed the coupled propagation and
decay treatment and quantified the dependence on the injected spectrum
\cite{Dutta:2000}; their later study showed that secondary muon neutrinos from
tau decay can contribute to upward-muon rates, with a particularly strong
spectral-shape dependence \cite{Dutta:2002}.  Rakshit and Reya studied
iterative solutions of coupled $\nu_\tau$--$\tau$ transport with secondary
$\nu_\mu$ production and convergence checks \cite{Rakshit:2006}.  These
works set the appropriate standard for a future topology-aware Case-6
propagation calculation.  The inherited v0.3 transport implementation was
checked against these historical benchmarks, but the present EHE fold does
not rerun that transport with the new differential kernel.

Modern Standard-Model neutrino DIS predictions include higher-order QCD,
small-$x$ PDF extrapolation and nuclear uncertainties.  Xie \emph{et al.}
provide CT18-based predictions over the UHE range and quantify those effects
\cite{Xie:2024}.  Our CT10nlo-based leading-order parton integral is a
controlled \emph{ratio} benchmark for the new operator; it is not a modern
precision Standard-Model cross-section calculation.  PDF, perturbative and
nuclear uncertainties must be revisited before a precision exclusion claim.

\subsection{The operator, collider information and cosmic flux}
Pandey, Karmakar and Rakshit studied monojet and neutrino constraints on a
range of non-standard interactions, including their Case 6, a derivative
dimension-seven neutrino--quark contact interaction.  Their quoted monojet
bound for its coefficient is
$|C_4|\lesssim 8.6\times10^{-10}\,\mathrm{GeV}^{-3}$
\cite{Pandey:2019}.  The earlier work compared cross sections and collider
limits, while the present audit combines its momentum-growing benchmark
with the subsequently released 12.6-year EHE response, the observed
energy-dependent limit and the specific ARCA bright-track topology.  We use
the quoted coefficient as a benchmark, not as an EFT-valid recast of LHC
events after partonic truncation.  A UV completion and a collider analysis
with the same validity cut would be required to make that correspondence
fully consistent.

Ehlert \emph{et al.} constrained high-energy proton fractions using
multimessenger cosmic-ray information and provided cosmogenic-neutrino
spectra that depend on source evolution and acceleration assumptions
\cite{Ehlert:2024}.  Choosing one of these shapes fixes the comparison, but
does not remove the astrophysical uncertainty in its normalization or
composition.  The KM3NeT open analysis notebook supplies a sky-averaged,
all-flavor bright-track effective area and per-flavor $\nu+\bar\nu$ flux
convention \cite{KM3NeT:2025Data}.  The full ARCA performance study treats
track and shower selections separately \cite{KM3NeT:2024ARCA}, which is
essential when the interaction has a neutral final state.

\subsection{Recent IceCube and KM3NeT comparisons}
IceCube's EHE analysis fits reconstructed energy and direction with a binned
Poisson likelihood, profiles an astrophysical component, and calibrates
tests with pseudoexperiments \cite{IceCube:2025EHE}.  Its release gives
true-energy effective areas and the published quasi-differential limit,
but not the full reconstructed-bin response \cite{IceCube:2025Data}.
KM3NeT's broader comparison of its event with the global neutrino landscape
uses joint flux fits \cite{KM3NeT:2025Landscape}.  A later analysis used
the single KM3NeT event to constrain the neutrino--nucleon cross section,
subject to its flux and detector hypotheses \cite{Bertolez:2026}; IceCube
also searched for point-source emission in the event direction
\cite{IceCube:2026PointSource}.  Those studies address complementary
questions.  Our specific novelty is the joint analytic flux-ratio and
topology audit of the Case-6 operator with both published public responses.
The result is conditional on a stationary isotropic common flux and cannot
exclude a transient or anisotropic source using this all-sky calculation.

\section{Analytic structure of Case 6}
The schematic low-energy interaction is
\begin{equation}
 {\cal L}_{\mathrm{C6}}= C_4\,
 (\bar\nu_i^{c}\overleftrightarrow{\partial}_{\mu}\nu_j)
 (\bar q\gamma^{\mu}q)+\mathrm{h.c.},\qquad [C_4]=-3.
 \label{eq:operator}
\end{equation}
The flavor and charge-conjugation assignments in a concrete theory determine
which incident species and regenerated species are affected.  The key
topology selection rule is robust: the hard scattering in
eq.~\eqref{eq:operator} produces a neutrino-sector final state and hadrons,
not a charged muon.  It adds no direct $\nu_\mu N\to\mu X$ track channel.
For an orthogonal neutral final state the new amplitude contributes a
non-negative additive term rather than interference with the Standard Model
charged-current track amplitude.

For a nucleon at rest and Bjorken variables $x,y$,
\begin{equation}
 \hat s\simeq 2m_NE_\nu x,\qquad Q^2=\hat s y,
 \qquad \overline{|{\cal M}_{\mathrm{C6}}|^2}
 =64 C_4^2\hat s^3y(1-y).
 \label{eq:kinematics}
\end{equation}
The polynomial vanishes at $y=0,1$ and peaks at $y=1/2$; the $\hat s^3$
factor explains the strong energy growth before PDF convolution.  The
primary EFT restriction is $\sqrt{\hat s}<\Lambda$, imposed point by
point.  The weaker $Q<\Lambda$ restriction admits configurations with
$Q^2=\hat s y<\Lambda^2$ while $\hat s>\Lambda^2$ and is reported only as an
aggressive envelope.  The dimensionless coefficient
$c_4=C_4\Lambda^3$ obeys $c_4/(4\pi)=0.231,0.547,0.994,1.848$ at the four
cutoffs used below.  Thus $\Lambda=3$ TeV crosses the simple $4\pi$
coefficient diagnostic, and $2.44$ TeV sits near it.  These diagnostics are
not a complete partial-wave unitarity analysis.  The expansion parameter
is $\hat s/\Lambda^2$; omitted higher-dimensional terms are uncontrolled
near the imposed boundary, so even the primary cut should be interpreted
as an EFT-domain envelope rather than a precision UV prediction.

Let $I$ and $K$ denote the predicted IceCube and KM3NeT counts per unit
common flux normalization, with the full energy, angle and topology
response understood.  Then the exact identity
\begin{equation}
 N_{\mathrm{IC}}(\kappa)=\kappa I,
 \quad N_{\mathrm{KM3}}(\kappa)=\kappa K,
 \quad \kappa\leq\frac{U}{I}
 \quad\Longrightarrow\quad
 N_{\mathrm{KM3}}\leq U\frac{K}{I}
 \label{eq:ratio}
\end{equation}
holds for any constraint of the indicated form.  If a proposed interaction
enhances $I$ and $K$ by similar factors, $K/I$ barely changes.  This
explains the near-cancellation seen in the initial v0.8 screen.  For an
actual track sample, Case 6 changes $K$ only indirectly through
propagation, flavor conversion or secondary-muon production; none is
quantified by simply multiplying its published track area by $1+r$.

\section{Public data and controlled screen}
For equal incident per-flavor fluxes $\phi_{\rm pf}$ and the IceCube release's
flavor-summed area, the Standard-Model event fold is
\begin{equation}
 N_{\rm IC}^{\rm SM}=4\pi T_{\rm IC}\int dE\,
 \phi_{\rm pf}(E)A_{\rm tot}^{\rm SM}(E).
 \label{eq:fold}
\end{equation}
The area is supplied in $\mathrm{m}^2$ and is multiplied by $10^4$ for use
with a flux in $\mathrm{cm}^{-2}$.  Its columns satisfy
$A_{\rm tot}\simeq A_e+A_\mu+A_\tau$ to tabulation precision.
The analogous KM3NeT fold uses its released all-flavor bright-track area
and $T_{\rm KM3}=335$ days.  We use $T_{\rm IC}=12.6$ Julian years.

To expose the topology uncertainty we compute two IceCube response
\emph{proxies}, never treating them as a re-simulation of the detector:
\begin{align}
 A_{\rm all}^{\rm proxy}(E)&=A_{\rm tot}^{\rm SM}(E)[1+r(E)],
 \\
 A_{\tau\text{-init}}^{\rm proxy}(E)&=
 A_{\rm tot}^{\rm SM}(E)+A_\tau^{\rm SM}(E)r(E),
 \qquad r=\sigma_{\rm C6}/\sigma_{\rm SM}.
 \label{eq:responseproxy}
\end{align}
These bracket choices about the incident flavors being affected, but they
do not bound the true detector response: neutral-current visible energy,
Earth attenuation and regeneration may change the selection in either
direction.  The published KM3NeT $A_{\rm eff}$ is for bright tracks.
We therefore keep the direct Case-6 contribution to that response exactly
zero.  Its hypothetical $A_{\rm KM3}^{\rm SM}(1+r)$ fold is retained in
the machine-readable table solely to show the size of a \emph{wrong-topology}
acceptance proxy.

IceCube's released differential-limit ordinate $L_j$ is an all-flavor
$E^2\Phi$ at central energy $E_j$.  The collaboration constructs each
point using an $E^{-1}$ test spectrum across a sliding one-decade window
\cite{IceCube:2025EHE}.  We reconstruct its true-energy reference flux as
\begin{equation}
 \phi_{j,{\rm pf}}^{\rm ref}(E)=
 \frac{L_j}{3E_j E}\,
 \Theta(E-E_j/\sqrt{10})\Theta(\sqrt{10}E_j-E),
 \qquad
 n_j^{\rm ref}=4\pi T_{\rm IC}\int dE\,
 A_{\rm tot}^{\rm SM}\phi_{j,{\rm pf}}^{\rm ref}.
 \label{eq:limitreference}
\end{equation}
For a chosen model response proxy, define $R_j=n_j^{\rm model}/n_j^{\rm ref}$
using the same window, and the diagnostic scale
\begin{equation}
 \kappa_{\rm window}^{\rm proxy}
 =\min_j R_j^{-1}.
 \label{eq:windowscale}
\end{equation}
This uses the \emph{observed} limit rather than the no-signal sensitivity.
Overlapping windows are correlated; the published limit is calibrated for
its $E^{-1}$ reference and Standard-Model response.  Equations
\eqref{eq:limitreference}--\eqref{eq:windowscale} do not inherit 90\%
coverage for a broad Ehlert spectrum or modified cross section.  The
published full likelihood would require event counts and signal/background
templates in reconstructed energy and direction, a migration response and
the nuisance treatment.  These are absent from the four-file public release.
Thus a genuine likelihood replacement cannot be produced from the released
tables alone.

\section{Results and interpretation}
The frozen v0.8 cross-section grid spans 0.1--20,000 PeV and uses CT10nlo
with the pointwise $\sqrt{\hat s}<\Lambda$ cut.  Its no-cut ratio
$r(0.1\,\mathrm{PeV})=5.34\times10^{-4}$ and
$r(1\,\mathrm{PeV})=1.38\times10^{-2}$ is preserved as a historical
regression.  The v0.9 numerical work only folds those saved curves with
the public data.  The KM3NeT Ehlert expectation is 0.242754 events, versus
0.242767 from its official notebook; in the published central 90\% energy
range it is 0.211081 versus 0.211117.  The corrected IceCube area fold is
17.6748 events at the unscaled Ehlert benchmark.  The difference from the
v0.8 value 17.6746 is the first numerical row of its headerless area file,
which the old parser silently skipped.  The original v0.8 unit error gave
0.001767 before the $\mathrm{m}^2\to\mathrm{cm}^2$ repair.

\begin{table}[t]
\centering
\caption{Version 0.9 public-data diagnostics.  $\kappa_{\rm all}$ and
$\kappa_{\tau}$ are the window-screen scales of
eq.~\eqref{eq:windowscale}, not confidence limits.  The last two columns
are the \emph{unchanged direct} bright-track expectation at those scales;
propagation effects have not been evaluated.}
\label{tab:results}
\begin{tabular}{@{}rrrrrrr@{}}
\toprule
$\Lambda$ [TeV] & $c_4/4\pi$ & $N_{\rm IC}^{\rm all}$ &
$N_{\rm IC}^{\tau}$ & $\kappa_{\rm all}$ &
$N_{\rm KM3}^{\rm tr}(\kappa_{\rm all})$ &
$N_{\rm KM3}^{\rm tr}(\kappa_{\tau})$\\
\midrule
1.50 & 0.231 & 20.734 & 18.885 & 0.1549 & 0.0376 & 0.0413\\
2.00 & 0.547 & 25.051 & 20.608 & 0.1280 & 0.0311 & 0.0378\\
2.44 & 0.994 & 31.151 & 23.049 & 0.1025 & 0.0249 & 0.0338\\
3.00 & 1.848 & 42.537 & 27.637 & 0.0747 & 0.0181 & 0.0281\\
\bottomrule
\end{tabular}
\end{table}

For all four benchmarks the most constraining released window is centered
at $E_j=3.162\times10^8$ GeV, or about $316$ PeV.  The all-flavor window
ratios there are 6.46, 7.81, 9.76 and 13.38, respectively.  The published
reference yield in that window is approximately 2.42 events when folded
with the released true-energy area.  These numbers show why the unscaled
cosmogenic benchmark is strongly disfavored by the public energy-dependent
screen.  They do not measure a frequentist model-rejection probability.

The reported $N_{\rm IC}^{\rm all}$ in table~\ref{tab:results} uses an
all-flavor $1+r$ acceptance proxy.  If only tau incident states are affected,
using $A_\tau$ reduces the effect as shown by $N_{\rm IC}^{\tau}$; the
corresponding $\kappa_\tau$ values are 0.1703, 0.1559, 0.1391 and 0.1157.
This flavor dependence is another reason that the v0.8 single-column proxy
cannot be interpreted as a physical event prediction.  In particular, the
v0.8 $N_{\rm KM3}^{\rm NC\ equivalent}$ enhancement is a diagnostic of
what happens if an NC cross section is assigned a track acceptance, not a
mechanism for the observed muon.  The direct track count at fixed incident
flux is the Standard-Model count in this calculation.

\section{Which alternatives remain viable?}
Equation~\eqref{eq:ratio} gives a necessary quantitative target for variants.
At $\Lambda=2$ TeV, the all-flavor screen leaves about 0.031 direct track
events, so a model seeking one expected bright track at the same screened
common flux must raise the \emph{track-to-IceCube yield ratio} by roughly
$1/0.031\simeq32$.  This is a necessary-condition diagnostic, not a
model-independent exclusion: the IceCube constraint and propagation must
be recomputed for each model.  The following distinctions matter.

An interaction with a charged-current muon final state can create tracks
directly, but the same interaction changes Earth opacity and IceCube's
tracks.  Its laboratory bounds, gauge embedding and EFT validity must be
derived before any detector comparison.  A neutral-current operator
coupled only to tau neutrinos can change tau regeneration or generate
secondary muon neutrinos through subsequent tau decays; this requires the
coupled transport kernels of \cite{Rakshit:2006,Dutta:2002} with the actual
Case-6 differential inelasticity distribution and flavor assignment.
Directional or time-dependent sources can break the common isotropic
exposure ratio; the present sky-averaged areas are inadequate for such
models.  Energy-localized or resonant new physics can also change the
IceCube/ARCA response ratio, but a contact EFT cannot be extrapolated
through its resonance, and the UV completion must obey laboratory and
unitarity constraints.  These variants are research paths, not positive
solutions demonstrated here.

\section{Reproducibility and limitations}
The project records v0.7 and v0.8 failed and corrected stages in
\texttt{PROGRESS.log}; v0.9 inputs and outputs are under versioned
\texttt{data/uhe\_official} and \texttt{results/v0.9}.  The one-dimensional
quadrature uses 6001 logarithmic energy points between $10^5$ and
$10^{11}$ GeV.  The released effective areas are linearly interpolated as
in the collaboration's benchmark notebook, and the tabulated Ehlert
$E^2\phi$ is interpolated in its provided logarithm.  The quadrature is
cheap; parallelizing it over 14 cores would add overhead and risk
oversubscription.  The expensive LHAPDF grid is reused rather than
recomputed.  Automated checks enforce area units, flavor-column sums,
and the two KM3NeT benchmark integrals.

The largest limitation is detector response under new physics.  The
published effective areas include selection and Standard-Model Earth
effects.  Multiplication by a cross-section ratio does not account for
neutral-current visible energy, energy migration, track/cascade
classification, flavor-changing propagation or angle dependence.  The
second limitation is statistical: the released differential limits are
not the model-specific binned likelihood.  Third, CT10nlo and leading-order
DIS are not precision inputs at all energies, and the $4\pi$ check is
only a perturbative diagnostic.  A definitive Case-6 test needs
collaboration-level response matrices or event simulation, a validated
modified Earth-transport calculation, and a consistently truncated
collider/UV interpretation.  Until then the v0.9 conclusion is that
momentum growth survives in the cross-section calculation, while its
naive bright-track explanation is unsupported and the common-flux
event-yield proxies are strongly constrained by IceCube.

\section*{AI-assisted drafting disclosure}
This working manuscript and its reproducibility scripts were prepared with
AI assistance under the author's direction.  Bibliographic records and
quoted public-data conventions were checked against the primary papers and
collaboration releases listed above.  The author should independently
verify every numerical and scientific claim before journal submission.

\bibliographystyle{JHEP}
\bibliography{references}
\end{document}
